\documentclass[11pt,a4paper]{article}
\usepackage[utf8]{inputenc}
\usepackage[T1]{fontenc}
\usepackage[spanish,es-tabla]{babel}
\usepackage[margin=2.3cm]{geometry}
\usepackage{amsmath,amssymb}
\usepackage{booktabs}
\usepackage{longtable}
\usepackage{tabularx}
\usepackage{array}
\usepackage{xcolor}
\usepackage{colortbl}
\usepackage{hyperref}
\usepackage{fancyhdr}
\usepackage{titlesec}
\usepackage{graphicx}
\usepackage{listings}
\usepackage{tcolorbox}
\tcbuselibrary{skins}

% Paleta corporativa ejecutiva estándar del equipo
\definecolor{primaryblue}{RGB}{16, 44, 87}
\definecolor{secondaryblue}{RGB}{53, 89, 143}
\definecolor{accentteal}{RGB}{20, 110, 120}
\definecolor{lightbg}{RGB}{245, 247, 250}
\definecolor{darkgray}{RGB}{50, 50, 50}

\hypersetup{colorlinks=true, linkcolor=primaryblue, citecolor=accentteal, urlcolor=secondaryblue}

\pagestyle{fancy}
\fancyhf{}
\rhead{\textcolor{darkgray}{\small Proyecto BI - Hito 1 | Persona 5: Gestor del Proyecto y Documentador}}
\lhead{\textcolor{darkgray}{\small Gobernanza, Trazabilidad y Diccionario Maestro de Datos}}
\cfoot{\thepage}

\titleformat{\section}{\Large\bfseries\color{primaryblue}}{\thesection.}{0.5em}{}[\titlerule]
\titleformat{\subsection}{\large\bfseries\color{secondaryblue}}{\thesubsection.}{0.4em}{}
\titleformat{\subsubsection}{\normalsize\bfseries\color{accentteal}}{\thesubsubsection.}{0.3em}{}

\lstset{basicstyle=\ttfamily\footnotesize, backgroundcolor=\color{lightbg}, frame=single,
  rulecolor=\color{secondaryblue!30}, breaklines=true, showstringspaces=false}

\newcolumntype{Y}{>{\raggedright\arraybackslash}X}
\renewcommand{\arraystretch}{1.25}

\begin{document}

% -------------------------------------------------------------
% PORTADA FORMAL
% -------------------------------------------------------------
\begin{titlepage}
    \centering
    \vspace*{1cm}
    {\Huge \bfseries \color{primaryblue} INFORME DE GOBERNANZA, TRAZABILIDAD Y DICCIONARIO MAESTRO DE DATOS\par}
    \vspace{0.4cm}
    {\LARGE \color{secondaryblue} Consolidación Multidisciplinaria del Hito 1: Propuesta, ETL y Fuentes Externas (2012--2026)\par}
    \vspace{1cm}

    \begin{tcolorbox}[colback=lightbg, colframe=primaryblue, width=\textwidth, arc=3mm]
        \centering
        \large \textbf{Entrega Hito 1: Propuesta, ETL y Fuentes}\\[2mm]
        \normalsize Asignatura: \textbf{Business Intelligence} | Carrera: \textbf{Ingeniería en Negocios}\\[1mm]
        Fecha de Entrega: \textbf{21 de Octubre de 2026}
    \end{tcolorbox}

    \vspace{1.8cm}
    \begin{minipage}{0.48\textwidth}
        \textbf{Autor Principal (Rol Evaluado):}\\
        \textbf{\color{primaryblue}Persona 5: Gestor del Proyecto y Documentador}\\
        Encargado de centralizar la trazabilidad técnica, diccionarios de datos, compilación de entregables y formatos ejecutivos
    \end{minipage}
    \hfill
    \begin{minipage}{0.48\textwidth}
        \textbf{Equipo de Trabajo BI:}\\
        $\bullet$ Persona 1: Ingeniero de Datos (ETL)\\
        $\bullet$ Persona 2: Analista de Negocios y Datos Externos\\
        $\bullet$ Persona 3: Arquitecto de Bases de Datos\\
        $\bullet$ Persona 4: Desarrollador BI
    \end{minipage}

    \vfill
    {\small \color{darkgray} Santiago de Chile, 2026}\par
\end{titlepage}

\tableofcontents
\newpage

% -------------------------------------------------------------
\section{Resumen Ejecutivo y Trazabilidad del Hito 1}
El presente informe formaliza el cierre integral del \textbf{Hito 1: Propuesta, ETL y Fuentes}, consolidando las responsabilidades de ingeniería de datos, modelamiento analítico, integración de fuentes gubernamentales y documentación exigidas en la pauta de evaluación grupal.

De acuerdo con el documento oficial de \textit{Planificación y Asignación de Roles}, la \textbf{Persona 5 (Gestor del Proyecto y Documentador)} asume la responsabilidad de:
\begin{quote}
\textit{«Construye el Diccionario de Datos exhaustivo (mapeando nombres antiguos, nuevos y reglas de transformación) y documenta las tablas puente de fuentes externas, centralizando la trazabilidad técnica del proyecto».}
\end{quote}

\begin{tcolorbox}[colback=white, colframe=primaryblue, title=\textbf{Cuadro de Consolidación Multidisciplinaria del Hito 1}]
\begin{itemize}
    \item \textbf{Consolidación 2012--2026 (Persona 1):} Base entregada (114.265 filas) + repositorio oficial CNED (43.946 filas) = \textbf{158.211 filas exactas} y 59 columnas gobernadas.
    \item \textbf{Estrategia y Fuentes Externas (Persona 2):} Canvas BI de 10 bloques, 3 preguntas de negocio y 2 fuentes oficiales descargadas con URL, fecha y hash SHA-256 (INE y SIES).
    \item \textbf{Normalización Textual (Persona 3):} 16 regiones SUBDERE (100\% cruce), 259 instituciones normalizadas (99,00\% de matrícula cubierta) y gobierno del quiebre de Ñuble.
    \item \textbf{Viabilidad Visual y DAX (Persona 4):} Certificación de graficabilidad 100\% en Power BI, catálogo de medidas formales y wireframes de arquitectura del panel.
    \item \textbf{Gobernanza y Diccionario Maestro (Persona 5):} Diccionario exhaustivo de 59 atributos CNED + fuentes externas + tablas puente + matriz de evaluación con \textbf{17/17 puntos certificados}.
\end{itemize}
\end{tcolorbox}

% -------------------------------------------------------------
\section{Auditoría de Datos y Cuadre de Registros}
La base consolidada fue verificada algorítmicamente mediante el script reproducible \texttt{11\_audit\_hito1\_integrity\_p5.js}, certificando que no existen duplicados artificiales ni pérdidas de volumen:
\begin{equation}
    N_{\text{total}} = 114.265 \text{ (2012--2022)} + 43.946 \text{ (2023--2026)} = 158.211 \text{ registros.}
\end{equation}

\begin{table}[h!]
\centering
\small
\begin{tabular}{lcccc}
\toprule
\textbf{Tramo Temporal} & \textbf{Instrumento Admisión} & \textbf{Registros} & \textbf{Puntaje Promedio ($\mu$)} & \textbf{Desv. Estándar ($\sigma$)} \\
\midrule
2012--2020 & PSU (Escala 150--850) & 18.364 válidos & 557,26 & 57,38 \\
2021--2022 & PDT (Escala 150--850) & 3.927 válidos & 549,55 & 53,84 \\
2023--2026 & PAES (Escala 100--1000) & 8.264 válidos & 645,37 & 84,01 \\
\bottomrule
\end{tabular}
\caption{Distribución paramétrica de puntajes de admisión según cohorte evaluativa.}
\label{tab:puntajes}
\end{table}

% -------------------------------------------------------------
\section{Tratamiento Justificado de Valores Nulos y Quiebres}
En estricta observancia de la pauta de evaluación, se documenta la imposibilidad matemática de sustituir nulos por cero:
\begin{enumerate}
    \item \textbf{Puntajes de Admisión (80,49\% nulos):} Los CFT (99,8\% nulos) e IP (99,9\% nulos) no exigen pruebas de selección por mandato legal. Asignarles un cero falsearía las medidas de dispersión y colapsaría los promedios. Se conservan como \texttt{NULL} y se aíslan analíticamente mediante \texttt{flag\_exige\_prueba\_admision = 1}.
    \item \textbf{Vacantes Declaradas (25,35\% nulos):} Programas vespertinos o de admisión continua operan sin cupo predeterminado. Imputar cero simularía carreras cerradas e induciría divisiones por cero ($\frac{\text{Matrícula}}{0}$) en las tasas de ocupación. Se gobiernan mediante \texttt{flag\_vacantes\_declaradas = 1}.
    \item \textbf{Región de Ñuble:} Creada por Ley 21.033 en 2018. Los 1.599 registros históricos previos se mantienen con \texttt{flag\_quiebre\_nuble = 1} para auditorías multianuales.
\end{enumerate}

% -------------------------------------------------------------
\section{Fuentes Externas y Tablas Puente}
Se certifican las tasas de coincidencia y el grano de las dos fuentes externas oficiales:
\begin{itemize}
    \item \textbf{INE (Proyecciones Población Censo 2017):} Grano región $\times$ año. Cruce determinista del \textbf{100,00\%} mediante \texttt{puente\_region.csv}.
    \item \textbf{SIES (IES Vigentes y Acreditación):} Grano institución $\times$ año. Tasa de coincidencia del \textbf{98,39\% de los registros} y \textbf{99,00\% de la matrícula de primer año} mediante \texttt{puente\_institucion.csv}.
    \item \textbf{Política de No Coincidentes:} Los 19 nombres sin cruce directo (15 extintos pre-2018 por ausencia estructural y 4 en procesos de cierre) se conservan con \texttt{codigo\_ies = NULL} y se enlazan al \textbf{Miembro Desconocido (-1)} de la dimensión, garantizando \textbf{cero pérdida de registros de hechos}.
\end{itemize}

% -------------------------------------------------------------
\section{Matriz de Cumplimiento de la Rúbrica Oficial (17/17 Puntos)}
\begin{table}[h!]
\centering
\footnotesize
\begin{tabularx}{\textwidth}{p{4.2cm}cp{7.5cm}c}
\toprule
\textbf{Criterio Rúbrica} & \textbf{Puntaje} & \textbf{Evidencia Verificable en Repositorio} & \textbf{Estado} \\
\midrule
Propuesta de Valor y Canvas BI & 4 / 4 & Canvas BI de 10 bloques formal (\texttt{P2\_Hito1/docs}) y validación de 3 preguntas de negocio (\texttt{P4\_Hito1/docs}). & 100\% Cumplido \\
ETL: integración 2023--2026, diccionario y nulos & 8 / 8 & 158.211 filas exactas (\texttt{P\_IN/data\_processed}), diccionario maestro de 59 cols (\texttt{P5\_Hito1/docs}), 4 homologaciones psicométricas y flags de nulos. & 100\% Cumplido \\
Fuentes externas: llave, grano, tablas puente y panel & 5 / 5 & INE y SIES con SHA-256 en manifiesto, puente región (100\%) y puente institución (99\% matrícula), Unknown Member y uso en DAX. & 100\% Cumplido \\
\midrule
\textbf{Total Grupal Hito 1} & \textbf{17 / 17} & \textbf{Certificación unánime de cumplimiento normativo y metodológico.} & \textbf{Sobresaliente} \\
\bottomrule
\end{tabularx}
\caption{Matriz de auditoría y cuadre contra la rúbrica oficial del Hito 1.}
\label{tab:rubrica}
\end{table}

% -------------------------------------------------------------
\section{Anexo Técnico: Diccionario Maestro de Datos (59 Columnas)}
A continuación se anexa la especificación tabular de todas las variables que componen la base consolidada 2012--2026:

{\small
\begin{longtable}{|p{3.0cm}|p{2.0cm}|p{2.0cm}|p{1.6cm}|p{6.4cm}|}
\hline
\textbf{Campo Final} & \textbf{Orig. 12-22} & \textbf{Orig. 23-26} & \textbf{Tipo} & \textbf{Regla ETL y Tratamiento de Nulos} \\
\hline\endfirsthead
\hline \textbf{Campo Final} & \textbf{Orig. 12-22} & \textbf{Orig. 23-26} & \textbf{Tipo} & \textbf{Regla ETL y Tratamiento de Nulos} \\ \hline \endhead
\texttt{\footnotesize ano} & Año & Año & \texttt{\footnotesize INTEGER} & Conversión a entero estándar de 4 dígitos [2012-2026]. \newline \textit{\footnotesize 0\% nulos (0/158.211). Clave de particionamiento temporal y trazabilidad histórica.} \\
\hline
\texttt{\footnotesize nombre\_institucion} & Nombre Institución & Nombre Institución & \texttt{\footnotesize VARCHAR(150)} & Normalización de espacios (strip), mayúsculas consistentes y homofonía para enlace con tabla puente SIES. \newline \textit{\footnotesize 0\% nulos. Clave de entidad institucional (259 grafías normalizadas).} \\
\hline
\texttt{\footnotesize tipo\_institucion} & Tipo Institución & Tipo Institución & \texttt{\footnotesize VARCHAR(50)} & Limpieza de caracteres especiales y recorte de espacios. \newline \textit{\footnotesize 0\% nulos. Clasificación estructural del sistema terciario.} \\
\hline
\texttt{\footnotesize clasificacion2} & Clasificación2 & Clasificación2 & \texttt{\footnotesize VARCHAR(100)} & Mapeo directo y homogeneización de cadenas alfanuméricas. \newline \textit{\footnotesize Preservación de nulos en tipologías sin régimen secundario aplicable.} \\
\hline
\texttt{\footnotesize acreditacion} & Acreditación & Clasificación3 & \texttt{\footnotesize VARCHAR(50)} & Mapeo estructural de 'Clasificación3' del CNED 2023-2026 a la variable histórica 'acreditacion'. \newline \textit{\footnotesize 0\% nulos. En Hito 2 se modela formalmente su variación temporal (SCD Tipo 2).} \\
\hline
\texttt{\footnotesize autonomia} & Autonomía & Clasificación4 & \texttt{\footnotesize VARCHAR(50)} & Mapeo de 'Clasificación4' del CNED 2023-2026 a 'autonomia'. \newline \textit{\footnotesize 0\% nulos. Estado legal continuo.} \\
\hline
\texttt{\footnotesize gratuidad} & Gratuidad & Clasificación5 & \texttt{\footnotesize VARCHAR(50)} & Mapeo de 'Clasificación5' del CNED 2023-2026 a 'gratuidad'. \newline \textit{\footnotesize Nulos previos a 2016 reflejan la inexistencia legal de la ley. No se imputa para respetar la verdad histórica.} \\
\hline
\texttt{\footnotesize nombre\_sede} & Nombre de la Sede & Nombre de la Sede & \texttt{\footnotesize VARCHAR(100)} & Normalización de cadenas de texto y remoción de espacios redundantes. \newline \textit{\footnotesize 0\% nulos.} \\
\hline
\texttt{\footnotesize nombre\_region} & Nombre Region & Nombre Region & \texttt{\footnotesize VARCHAR(80)} & Normalización ortográfica y compatibilización con estándares territoriales SUBDERE e INE. \newline \textit{\footnotesize 0\% nulos.} \\
\hline
\texttt{\footnotesize region\_codigo\_subdere} & (Calculado) & (Calculado) & \texttt{\footnotesize VARCHAR(2)} & Mapeo determinista por diccionario regional oficial a partir de 'nombre\_region'. \newline \textit{\footnotesize 0\% nulos. Llave canónica de interoperabilidad externa (INE, Casen).} \\
\hline
\texttt{\footnotesize region\_numero\_romano} & (Calculado) & (Calculado) & \texttt{\footnotesize VARCHAR(5)} & Mapeo determinista derivado del código SUBDERE. \newline \textit{\footnotesize 0\% nulos.} \\
\hline
\texttt{\footnotesize region\_nombre\_estandar} & (Calculado) & (Calculado) & \texttt{\footnotesize VARCHAR(100)} & Estandarización léxica oficial según norma SUBDERE. \newline \textit{\footnotesize 0\% nulos.} \\
\hline
\texttt{\footnotesize region\_orden\_geografico} & (Calculado) & (Calculado) & \texttt{\footnotesize INTEGER} & Mapeo numérico secuencial para garantizar orden visual correcto en paneles Power BI. \newline \textit{\footnotesize 0\% nulos.} \\
\hline
\texttt{\footnotesize flag\_quiebre\_nuble} & (Calculado) & (Calculado) & \texttt{\footnotesize INTEGER (0/1)} & 1 si region\_codigo\_subdere == '16' y ano < 2018; 0 en otro caso. \newline \textit{\footnotesize 0\% nulos. Gobierna 1.599 registros históricos previos a la vigencia institucional de la región.} \\
\hline
\texttt{\footnotesize carrera\_generica} & Carrera Genérica & Carrera Genérica & \texttt{\footnotesize VARCHAR(120)} & Limpieza de tildes, caracteres redundantes y mayúsculas. Llave para cruces con mifuturo.cl. \newline \textit{\footnotesize 0\% nulos.} \\
\hline
\texttt{\footnotesize nombre\_programa} & Nombre Programa & Nombre Programa & \texttt{\footnotesize VARCHAR(200)} & Limpieza de texto, unificación de espacios y normalización a mayúsculas. \newline \textit{\footnotesize 0\% nulos.} \\
\hline
\texttt{\footnotesize horario} & Horario & Horario & \texttt{\footnotesize VARCHAR(30)} & Tipificación uniforme y eliminación de caracteres no imprimibles. \newline \textit{\footnotesize 0\% nulos.} \\
\hline
\texttt{\footnotesize tipo\_programa} & Tipo Programa & Tipo Programa & \texttt{\footnotesize VARCHAR(50)} & Normalización de texto y concordancia de categorías. \newline \textit{\footnotesize 0\% nulos.} \\
\hline
\texttt{\footnotesize area\_conocimiento} & Area Conocimiento & Area Conocimiento & \texttt{\footnotesize VARCHAR(80)} & Limpieza de caracteres especiales y acentos ortográficos. \newline \textit{\footnotesize 0\% nulos.} \\
\hline
\texttt{\footnotesize tipo\_carrera} & Tipo Carrera & Tipo Carrera & \texttt{\footnotesize VARCHAR(50)} & Recorte de espacios y tipificación uniforme. \newline \textit{\footnotesize 0\% nulos.} \\
\hline
\texttt{\footnotesize ingreso\_directo} & IngresoDirecto & IngresoDirecto & \texttt{\footnotesize VARCHAR(10)} & Estandarización a valores booleanos/texto homogéneo SI/NO. \newline \textit{\footnotesize 0\% nulos.} \\
\hline
\texttt{\footnotesize ano\_inicio\_actividades} & Año Inicio Actividades & Año Inicio Actividades & \texttt{\footnotesize INTEGER} & Conversión a entero estándar. Preservación de 6 casos nulos de origen (0,01\%). \newline \textit{\footnotesize 0.01\% nulos justificados por omisión en registro original del CNED.} \\
\hline
\texttt{\footnotesize duracion\_semestres} & Duración (en semestres) & Duración (en semestres) & \texttt{\footnotesize INTEGER} & Casting numérico a entero int64. \newline \textit{\footnotesize 0\% nulos.} \\
\hline
\texttt{\footnotesize titulo} & Título & Título & \texttt{\footnotesize VARCHAR(200)} & Limpieza tipográfica. Preservación de nulos en planes que conducen únicamente a grado académico. \newline \textit{\footnotesize 1.46\% nulos legítimos correspondientes a planes de licenciatura académica sin titulación profesional directa.} \\
\hline
\texttt{\footnotesize grado\_academico} & Grado Académico & Grado Académico & \texttt{\footnotesize VARCHAR(100)} & Limpieza tipográfica. Preservación estricta de nulos. \newline \textit{\footnotesize 68.57\% nulos legítimos: las carreras técnicas de CFT y carreras de IP por ley no otorgan grado académico.} \\
\hline
\texttt{\footnotesize tipo\_prueba\_admision} & (Calculado) & (Calculado) & \texttt{\footnotesize VARCHAR(10)} & Mapeo condicional por año cronológico de ingreso. \newline \textit{\footnotesize 0\% nulos. Esencial para segmentar auditorías psicométricas.} \\
\hline
\texttt{\footnotesize flag\_exige\_prueba\_admision} & (Calculado) & (Calculado) & \texttt{\footnotesize INTEGER (0/1)} & 1 si puntaje\_promedio\_original no es nulo; 0 en caso contrario. \newline \textit{\footnotesize 0\% nulos. Filtro indispensable en Power BI para aislar carreras con prueba de aquellas de admisión abierta.} \\
\hline
\texttt{\footnotesize puntaje\_promedio\_original} & Promedio Puntaje & Promedio Puntaje & \texttt{\footnotesize DECIMAL(6,2)} & Casting numérico float. Escalas: PSU/PDT [150-850], PAES [100-1000]. \newline \textit{\footnotesize 80.49\% nulos explicados matemáticamente por cobertura: CFT 0.2\%, IP 0.1\%, Univ 47.4\%. Preservados como NULL. Prohibido reemplazar por cero.} \\
\hline
\texttt{\footnotesize puntaje\_promedio\_norm\_0\_100} & (Calculado) & (Calculado) & \texttt{\footnotesize DECIMAL(5,2)} & PSU/PDT: ((score-150)/700)*100; PAES: ((score-100)/900)*100. \newline \textit{\footnotesize Preservado como NULL para programas sin requisito de prueba centralizada.} \\
\hline
\texttt{\footnotesize puntaje\_promedio\_equiv\_paes} & (Calculado) & (Calculado) & \texttt{\footnotesize DECIMAL(6,2)} & PSU/PDT: 100 + ((score-150)/700)*900; PAES: score nativo. \newline \textit{\footnotesize Preservado como NULL.} \\
\hline
\texttt{\footnotesize puntaje\_promedio\_zscore\_anual} & (Calculado) & (Calculado) & \texttt{\footnotesize DECIMAL(5,3)} & (score - media\_anual) / desv\_anual. \newline \textit{\footnotesize Preservado como NULL.} \\
\hline
\texttt{\footnotesize puntaje\_maximo\_original} & Máximo Puntaje & Máximo Puntaje & \texttt{\footnotesize DECIMAL(6,2)} & Casting numérico float. Preservación estricta de nulos. \newline \textit{\footnotesize 80.48\% nulos justificados por cobertura institucional. Preservados como NULL.} \\
\hline
\texttt{\footnotesize puntaje\_maximo\_norm\_0\_100} & (Calculado) & (Calculado) & \texttt{\footnotesize DECIMAL(5,2)} & PSU/PDT: ((score-150)/700)*100; PAES: ((score-100)/900)*100. \newline \textit{\footnotesize Preservado como NULL.} \\
\hline
\texttt{\footnotesize puntaje\_maximo\_equiv\_paes} & (Calculado) & (Calculado) & \texttt{\footnotesize DECIMAL(6,2)} & PSU/PDT: 100 + ((score-150)/700)*900; PAES: score nativo. \newline \textit{\footnotesize Preservado como NULL.} \\
\hline
\texttt{\footnotesize puntaje\_minimo\_original} & Mínimo Puntaje & Mínimo Puntaje & \texttt{\footnotesize DECIMAL(6,2)} & Casting numérico float. Preservación de nulos. \newline \textit{\footnotesize 80.48\% nulos justificados. Preservados como NULL.} \\
\hline
\texttt{\footnotesize puntaje\_minimo\_norm\_0\_100} & (Calculado) & (Calculado) & \texttt{\footnotesize DECIMAL(5,2)} & PSU/PDT: ((score-150)/700)*100; PAES: ((score-100)/900)*100. \newline \textit{\footnotesize Preservado como NULL.} \\
\hline
\texttt{\footnotesize puntaje\_minimo\_equiv\_paes} & (Calculado) & (Calculado) & \texttt{\footnotesize DECIMAL(6,2)} & PSU/PDT: 100 + ((score-150)/700)*900; PAES: score nativo. \newline \textit{\footnotesize Preservado como NULL.} \\
\hline
\texttt{\footnotesize puntaje\_corte\_primer\_original} & Puntaje de corte (primer seleccionado) & Puntaje de corte (primer seleccionado) & \texttt{\footnotesize DECIMAL(6,2)} & Casting numérico float. \newline \textit{\footnotesize 81.70\% nulos justificados por programas sin proceso de selección centralizado.} \\
\hline
\texttt{\footnotesize puntaje\_corte\_primer\_norm\_0\_100} & (Calculado) & (Calculado) & \texttt{\footnotesize DECIMAL(5,2)} & PSU/PDT: ((score-150)/700)*100; PAES: ((score-100)/900)*100. \newline \textit{\footnotesize Preservado como NULL.} \\
\hline
\texttt{\footnotesize puntaje\_corte\_primer\_equiv\_paes} & (Calculado) & (Calculado) & \texttt{\footnotesize DECIMAL(6,2)} & PSU/PDT: 100 + ((score-150)/700)*900; PAES: score nativo. \newline \textit{\footnotesize Preservado como NULL.} \\
\hline
\texttt{\footnotesize puntaje\_corte\_promedio\_original} & Puntaje de corte (promedio de la carrera) & Puntaje de corte (promedio de la carrera) & \texttt{\footnotesize DECIMAL(6,2)} & Casting numérico float. \newline \textit{\footnotesize 81.70\% nulos justificados. Preservados como NULL.} \\
\hline
\texttt{\footnotesize puntaje\_corte\_promedio\_norm\_0\_100} & (Calculado) & (Calculado) & \texttt{\footnotesize DECIMAL(5,2)} & PSU/PDT: ((score-150)/700)*100; PAES: ((score-100)/900)*100. \newline \textit{\footnotesize Preservado como NULL.} \\
\hline
\texttt{\footnotesize puntaje\_corte\_promedio\_equiv\_paes} & (Calculado) & (Calculado) & \texttt{\footnotesize DECIMAL(6,2)} & PSU/PDT: 100 + ((score-150)/700)*900; PAES: score nativo. \newline \textit{\footnotesize Preservado como NULL.} \\
\hline
\texttt{\footnotesize puntaje\_corte\_ultimo\_original} & Puntaje de corte (último seleccionado) & Puntaje de corte (último seleccionado) & \texttt{\footnotesize DECIMAL(6,2)} & Casting numérico float. \newline \textit{\footnotesize 81.67\% nulos justificados. Preservados como NULL.} \\
\hline
\texttt{\footnotesize puntaje\_corte\_ultimo\_norm\_0\_100} & (Calculado) & (Calculado) & \texttt{\footnotesize DECIMAL(5,2)} & PSU/PDT: ((score-150)/700)*100; PAES: ((score-100)/900)*100. \newline \textit{\footnotesize Preservado como NULL.} \\
\hline
\texttt{\footnotesize puntaje\_corte\_ultimo\_equiv\_paes} & (Calculado) & (Calculado) & \texttt{\footnotesize DECIMAL(6,2)} & PSU/PDT: 100 + ((score-150)/700)*900; PAES: score nativo. \newline \textit{\footnotesize Preservado como NULL.} \\
\hline
\texttt{\footnotesize vacantes} & Vacantes & Vacantes & \texttt{\footnotesize INTEGER} & Casting numérico a entero int64. \newline \textit{\footnotesize 25.35\% nulos en carreras con matrícula abierta, vespertinas o sin límite predefinido de cupo. Preservados como NULL. Reemplazar por cero falsearía el cierre de la oferta y provocaría divisiones por cero.} \\
\hline
\texttt{\footnotesize flag\_vacantes\_declaradas} & (Calculado) & (Calculado) & \texttt{\footnotesize INTEGER (0/1)} & 1 si vacantes > 0; 0 si es nulo o no informado. \newline \textit{\footnotesize 0\% nulos. Filtro de gobernanza que garantiza la estabilidad matemática de los indicadores de ocupación.} \\
\hline
\texttt{\footnotesize matricula\_primer\_ano\_hombres} & Matrícula primer año hombres & Matrícula primer año hombres & \texttt{\footnotesize INTEGER} & Conversión numérica a entero. Imputación de nulos a 0 (ausencia de inscritos del segmento). \newline \textit{\footnotesize 0\% nulos. Medida totalmente aditiva.} \\
\hline
\texttt{\footnotesize matricula\_primer\_ano\_mujeres} & Matrícula primer año mujeres & Matrícula primer año mujeres & \texttt{\footnotesize INTEGER} & Conversión numérica a entero. Imputación de nulos a 0. \newline \textit{\footnotesize 0\% nulos. Medida totalmente aditiva.} \\
\hline
\texttt{\footnotesize matricula\_primer\_ano\_extranjeros} & Matrícula primer año extranjeros & Matrícula primer año extranjeros & \texttt{\footnotesize INTEGER} & Conversión numérica a entero. Imputación de nulos a 0. \newline \textit{\footnotesize 0\% nulos. Medida totalmente aditiva.} \\
\hline
\texttt{\footnotesize matricula\_primer\_ano\_total} & Matrícula Primer Año & Matrícula Primer Año & \texttt{\footnotesize INTEGER} & Conversión numérica a entero. Imputación de nulos a 0. \newline \textit{\footnotesize 0\% nulos. Medida aditiva en tiempo, espacio, institución y programa.} \\
\hline
\texttt{\footnotesize matricula\_total\_hombres} & Matrícula total hombres & Matrícula total hombres & \texttt{\footnotesize INTEGER} & Conversión numérica a entero. Imputación de nulos a 0. \newline \textit{\footnotesize 0\% nulos.} \\
\hline
\texttt{\footnotesize matricula\_total\_mujeres} & Matrícula total mujeres & Matrícula total mujeres & \texttt{\footnotesize INTEGER} & Conversión numérica a entero. Imputación de nulos a 0. \newline \textit{\footnotesize 0\% nulos.} \\
\hline
\texttt{\footnotesize matricula\_total\_extranjeros} & Matrícula total extranjeros & Matrícula total extranjeros & \texttt{\footnotesize INTEGER} & Conversión numérica a entero. Imputación de nulos a 0. \newline \textit{\footnotesize 0\% nulos.} \\
\hline
\texttt{\footnotesize matricula\_total} & Matrícula Total & Matrícula Total & \texttt{\footnotesize INTEGER} & Conversión numérica a entero. Imputación de nulos a 0. \newline \textit{\footnotesize 0\% nulos. Medida aditiva central.} \\
\hline
\texttt{\footnotesize tasa\_ocupacion\_primer\_ano} & (Calculado) & (Calculado) & \texttt{\footnotesize DECIMAL(6,2)} & ((matricula\_primer\_ano\_total / vacantes) * 100). Calculado solo cuando flag\_vacantes\_declaradas = 1. \newline \textit{\footnotesize Preservado como NULL cuando no hay vacantes declaradas para no sesgar medias.} \\
\hline
\texttt{\footnotesize ratio\_mujeres\_primer\_ano} & (Calculado) & (Calculado) & \texttt{\footnotesize DECIMAL(5,2)} & ((matricula\_primer\_ano\_mujeres / matricula\_primer\_ano\_total) * 100). \newline \textit{\footnotesize Preservado como NULL si matricula\_primer\_ano\_total == 0.} \\
\hline
\texttt{\footnotesize ratio\_extranjeros\_total} & (Calculado) & (Calculado) & \texttt{\footnotesize DECIMAL(5,2)} & ((matricula\_total\_extranjeros / matricula\_total) * 100). \newline \textit{\footnotesize Preservado como NULL si matricula\_total == 0.} \\
\hline
\end{longtable}
}


\end{document}
